\chapter{The big picture}
\label{ch:bigpicture}

\chaptergoals{2 of 11 in Part~II.}
{\cref{ch:setup}: a working environment.}
{name the packages of the chain and the file each one produces, read a Phlex workflow, and explain why runs are reproducible.}

\section{The guiding example}

A heavy neutral lepton mixes with the active neutrinos. At SHiP it is produced
mostly in decays of charm and beauty mesons in the target, flies through the
muon shield without interacting, and decays somewhere in the 50\,m long decay
volume. We follow its decay
\[
  \HNL \to \pi^{\pm}\mu^{\mp},
\]
a useful benchmark signature: two charged tracks from a common vertex
inside the decay volume, no activity in the veto taggers, and a reconstructed invariant
mass distributed around the HNL mass.

You will take a sample of such decays and push it through the whole chain in
\cref{fig:chain}: transport the two daughters through the detector, look at
the result in the event display, turn energy deposits into detector hits, and
finally form an illustrative weighted acceptance estimate.

\begin{figure}[h]
\centering
\begin{tikzpicture}[
  node distance=7mm and 6mm,
  every node/.style={font=\sffamily\small},
  stage/.style={draw=shipnavy, fill=shipcream!60, rounded corners=2pt,
    minimum width=27mm, minimum height=9mm, align=center, line width=0.7pt},
  service/.style={draw=shipblue, fill=shipblue!8, rounded corners=2pt,
    minimum width=27mm, minimum height=8mm, align=center},
  datafile/.style={draw=shippink, fill=shippink!8, rounded corners=0pt,
    minimum width=23mm, minimum height=6mm,
    align=center, font=\ttfamily\footnotesize},
  future/.style={stage, draw=black!65, fill=gray!8, densely dotted, text=black!75},
  flow/.style={-{Stealth[length=2.2mm]}, line width=0.8pt, color=shipnavy},
  uses/.style={-{Stealth[length=2mm]}, color=shipblue, dashed},
]
  \node[stage] (evc) {EventCalc\\\scriptsize HNL decays};
  \node[datafile, below=of evc] (hnl) {hnl\_events.root};
  \node[stage, below=of hnl] (aegir) {\textbf{aegir}\\\scriptsize Geant4 transport};
  \node[datafile, below=of aegir] (sim) {hnl\_sim.root};
  \node[stage, below=of sim] (shannon) {\textbf{Shannon}\\\scriptsize digitisation};
  \node[datafile, below=of shannon] (digi) {hnl\_digi.root};
  \node[future, below=of digi] (trout) {Trout\\\scriptsize reconstruction};

  \node[service, left=12mm of aegir] (geosvc) {geometry service};
  \node[service, above=of geosvc] (geo) {Geometry\\\scriptsize\ttfamily ship\_geometry.db};
  \node[service, below=of geosvc] (field) {field service\\\scriptsize field maps};

  \node[stage, right=12mm of sim] (sc) {\textbf{sea-cucumber}\\\scriptsize event display\\[-1pt]\scriptsize(+ \texttt{ship\_geometry.db})};
  \node[stage, right=12mm of digi, fill=shipcream!25] (ana) {your analysis\\\scriptsize Python + ROOT};

  \draw[flow] (evc) -- node[right, font=\scriptsize\sffamily] {your converter} (hnl);
  \draw[flow] (hnl) -- (aegir);
  \draw[flow] (aegir) -- (sim);
  \draw[flow] (sim) -- (shannon);
  \draw[flow] (shannon) -- (digi);
  \draw[flow, black!65, densely dotted] (digi) -- (trout);

  \draw[uses] (geo) -- (geosvc);
  \draw[uses] (geosvc) -- (aegir);
  \draw[uses] (field) -- (aegir);
  \draw[flow] (sim) -- (sc);
  \draw[flow] ([yshift=-2mm]sim.east) -- ++(5mm,0) |- ([yshift=2mm]ana.west);
  \draw[flow] (digi) -- (ana);

  \begin{scope}[on background layer]
    \node[draw=shipblue!40, fill=shipblue!3, rounded corners=4pt, inner sep=4mm,
      fit=(geo)(geosvc)(field), label={[font=\sffamily\scriptsize\color{shipblue}]above:detector description}] {};
  \end{scope}
\end{tikzpicture}
\caption{The chain followed in this booklet. Solid arrows carry event data
between stages; dashed arrows supply detector information. The dotted
Trout branch is an outlook. Event files are ROOT RNTuples; the geometry
is a SQLite database and field maps are ROOT RNTuples in their own schema.}
\label{fig:chain}
\end{figure}

\section{The packages}

\begin{center}
\small
\begin{tabularx}{\linewidth}{l l L}
\toprule
Package & Repository & Role \\
\midrule
\pkg{Geometry}          & \repo{Geometry}         & The SHiP detector described in GeoModel; writes \file{ship\_geometry.db}. \\
\pkg{geometry service}  & \repo{geometry\_service}& Framework-independent access to that geometry, as GeoModel or as Geant4 volumes. \\
\pkg{field service}     & \repo{field\_service}   & Field maps in a versioned file format, evaluated for Geant4 and ACTS. \\
\pkg{data model}        & \repo{data-model}       & The event classes (\texttt{MCParticle}, \texttt{SimHit}, \ldots) and their ROOT I/O. \\
\pkg{ship-random}       & \repo{ship-random}      & The counter-based random number generator shared by all packages. \\
\pkg{aegir}             & \repo{aegir}            & Simulation: event sources plus Geant4, as Phlex plugins. \\
\pkg{sea-cucumber}      & \repo{sea-cucumber}     & Event display, in ROOT REve or in the browser. \\
\pkg{Shannon}           & \repo{Shannon}          & Digitisation: simulated hits into detector hits. \\
\pkg{Trout}             & \repo{Trout}            & Reconstruction (not covered hands-on, see \cref{ch:outlook}). \\
\bottomrule
\end{tabularx}
\end{center}

All of these are new code. \pkg{FairShip}, the framework used for SHiP so far,
is not needed for anything in this booklet.

\section{Phlex in one page}
\label{sec:phlex}

\pkg{aegir}, \pkg{Shannon} and \pkg{Trout} are not programs in their own right:
they are collections of plugins for
\href{https://github.com/Framework-R-D/phlex}{Phlex}, a data-flow framework
developed by the Framework R\&D project. You never write a \texttt{main()}. Instead you write a
\emph{workflow} file that names the plugins to load and configures each one,
and the \cmd{phlex} executable builds a graph from it and runs it, in
parallel where the data allow.

\paragraph{Plugins.} Each plugin is a shared library found on
\texttt{PHLEX\_PLUGIN\_PATH}, referred to in a workflow by its \texttt{cpp}
name. There are four kinds you will meet:

\begin{description}[leftmargin=2.2cm, style=sameline]
  \item[driver]   says which data cells exist: for example ``a job with 100
                  events'' (\texttt{generate\_layers}).
  \item[source / provider] puts products into a cell: generated or read
                  particles per event, the geometry and field once per job.
  \item[transform] turns products into new products, such as particles into
                  hits (Geant4) or hits into digitised hits.
  \item[observer] consumes products without producing any: output writers and
                  histogram fillers.
\end{description}

\paragraph{Layers.} Cells are organised in layers: a \texttt{job} contains
\texttt{event}s in \pkg{aegir}, and \texttt{spill}s in \pkg{Shannon}. A product
lives in one layer; the geometry is a job-level product shared by every event.

\paragraph{Stages.} Every job names its processing stage
(\texttt{simulation}, \texttt{digitisation}), and the products it creates are
labelled with it.

\paragraph{Workflows are jsonnet.}
\href{https://jsonnet.org}{Jsonnet} is JSON with variables, functions and
imports. The \pkg{aegir} workflows import building blocks from
\file{workflows/lib.libsonnet} and override what they need. Values that change
from run to run (number of events, file names) are passed as external
variables, so a run is typically:

\begin{shell}
phlex -c <(jsonnet --ext-str events=100 --ext-str infile=hnl_events.root \
                   my_workflow.jsonnet) -j 4
\end{shell}

\cmd{jsonnet} expands the workflow to plain JSON, the shell hands it to
\cmd{phlex} as a file, and \cmd{-j 4} lets Phlex use four threads.

\begin{note}
You can always look at what \cmd{phlex} actually receives by running the
\cmd{jsonnet} part on its own. This is the first thing to do when a workflow
does not behave as expected.
\end{note}

\section{Reproducibility in one paragraph}

The simulation and digitisation plugins use \pkg{ship-random}'s
\texttt{PhiloxRng}, a counter-based generator: a fresh generator is seeded per
event from \emph{(seed, stream key, event number)}, with no state shared
between events. The same seed therefore gives the same events whatever the
number of threads and in whatever order events finish. The aegir generator and transport plugins choose and log a random seed
when none is supplied; Shannon defaults to seed zero. This booklet sets
explicit seeds. Reproducibility also requires fixed inputs, software and
configuration. Shannon seeds by input entry/spill index, so reordering a
simulation file changes which random stream a physical event receives.

\begin{pitfall}
Output writers fill RNTuples from several threads at once, so entries appear
in the order events \emph{finish}, not in event-number order. In \pkg{aegir}
output every entry carries its \field{event\_number}; never assume that entry
$i$ is event $i$.
\end{pitfall}
